\documentclass[11pt]{article}
\usepackage[T1]{fontenc}
\usepackage[utf8]{inputenc}
\usepackage{lmodern,microtype}
\usepackage[margin=0.95in]{geometry}
\usepackage{amsmath,amssymb,amsthm,mathtools,booktabs,graphicx}
\usepackage{xcolor}
\usepackage[colorlinks=true,linkcolor=blue!55!black,citecolor=blue!55!black]{hyperref}
\setlength{\emergencystretch}{3em}
\newtheorem{theorem}{Theorem}[section]
\newtheorem{lemma}[theorem]{Lemma}
\newtheorem{proposition}[theorem]{Proposition}
\newtheorem{corollary}[theorem]{Corollary}
\newcommand{\R}{\mathbb R}
\newcommand{\E}{\mathbb E}
\newcommand{\N}{\mathcal N}
\newcommand{\norm}[1]{\left\lVert #1\right\rVert}
\newcommand{\clip}{\operatorname{clip}}
\newcommand{\sg}{\mathrm{sg}}
\newcommand{\gap}{\mathrm{gap}}
\title{DP-SGD batch-size theory with the analytical GDP approximation\\
\large Convergence, optimal learning rates, and a separate finite privacy guarantee}
\author{}
\date{}
\begin{document}
\maketitle
\begin{abstract}
We study how privacy calibration affects batch-size predictions for
Poisson-sampled, per-example clipped SGD. The analytical GDP approximation
provides a logarithmic noise prescription, retaining batch dependence that
cancels in the Abadi scaling prescription. For any specified positive noise
variance, we prove a finite bound on the average expected gradient norm and
the learning rate minimizing that bound. We then specialize these exact
optimization results to the GDP approximation, distinguish fixed updates
from fixed expected passes, and compare with numerical privacy accounting.
At fixed expected passes, we show that the inverse-CLT prescription has a
degenerate privacy limit, so its batch-size predictions require a model-based
interpretation. An explicit finite event also disproves a universal transcript
privacy guarantee at this noise. A separate appendix gives a sufficient
finite calibration. Numerical accounting supplies the practical comparator.
\end{abstract}

\section{Algorithm and main results}
\label{sec:main}
The utility proof adapts the clipping regimes and descent argument of
Koloskova, Hendrikx and Stich~\cite[Appendices C.3--C.4]{clipping}, with
explicit Poisson sampling moments and constants. Selecting a learning rate
by minimizing a convergence bound also underlies Section~4.2 of Monzio
Compagnoni et al.~\cite{amap}; we derive the optimizer of the bound given here.
For the main batch-size model we use the analytical GDP approximation of
Bu et al.~\cite{bu2020deepgdp}, based on Gaussian differential
privacy~\cite{gdp}. Its privacy interpretation is asymptotic; the convergence
bound below is finite and holds for every chosen $V>0$.
Appendix~\ref{sec:privacy} provides a separate sufficient finite privacy
calibration. Neither that sufficient calibration nor the CLT approximation
is asserted to equal the smallest feasible noise variance.

\paragraph{Sampling and adjacency.}
Fix public integers $N,d,T\ge1$, a sampling probability $q\in(0,1]$, a
clipping threshold $C>0$, and an initialization $X_0\in\R^d$.
A dataset consists of $N$ slots, each empty or containing one record.
Neighboring datasets differ by the addition or removal of one record in one
slot. The public $N$ and the denominator $B=qN$ stay fixed across neighbors.
A common record-query function supplies a measurable vector $h(r,x)$;
an empty slot contributes zero. Write $h_i(x)$ for the query in slot $i$ and
let $\clip_C$ denote Euclidean projection onto the ball of radius $C$.
For fixed public $\eta>0$ and multiplier variance $V>0$, use
\begin{equation}\label{eq:algorithm}
 Y_t=\sum_{i=1}^N B_{t,i}\clip_C(h_i(X_t))+C\sqrt V\,Z_t,
 \qquad X_{t+1}=X_t-\frac{\eta}{B}Y_t,
 \quad 0\le t<T.
\end{equation}
Here all $B_{t,i}\sim\operatorname{Bernoulli}(q)$ and $Z_t\sim\N(0,I_d)$
are independent. Sampling masks and noise draws remain internal. The
privacy theorem protects the entire noisy-sum transcript $(Y_t)_{t<T}$,
and hence any public postprocessing such as the final iterate.

\paragraph{Noise units.}
The noise multiplier is $\sigma_{\mathrm{mult}}=\sqrt V$.
After division by $B$, the Gaussian noise has coordinate standard deviation
$C\sqrt V/B$ and total noise second moment
\begin{equation}\label{eq:noise}
 s^2=\E\norm{C\sqrt V Z_t/B}^2=\frac{C^2dV}{B^2}.
\end{equation}
The stepsize multiplies this noise subsequently. The quantity $s$ is
distinct from the stochastic-gradient standard deviation $\sigma_{\sg}$.

\paragraph{Analytical GDP calibration.}
Let $\Phi$ and $\varphi$ denote the standard normal distribution function and
density. For separation $\mu>0$, the exact Gaussian privacy profile is
\begin{equation}\label{eq:gaussian-profile}
 g_\mu(z)=\Phi\left(\frac\mu2-\frac z\mu\right)
       -e^z\Phi\left(-\frac\mu2-\frac z\mu\right).
\end{equation}
For public $\varepsilon>0$ and $\delta\in(0,1)$, select the unique
$\mu>0$ satisfying $g_\mu(\varepsilon)=\delta$. This scalar inversion is
numerical in general: $\partial_\mu g_\mu(\varepsilon)
=\varphi(\mu/2-\varepsilon/\mu)>0$, and the endpoint limits are zero and one.
The subsampled CLT model and its inverse are
\begin{equation}\label{eq:clt-calibration}
 \mu_{\mathrm{CLT}}^2=Tq^2(e^{1/V}-1),\qquad
 V_{\mathrm{CLT}}(T,q;\mu)
 =\frac{1}{\log(1+\mu^2/(Tq^2))},\quad 0<q<1.
\end{equation}
Thus the variance is closed form in $(T,q,\mu)$; the full calibration from
$(\varepsilon,\delta)$ also involves the one-dimensional Gaussian inversion.
At $q=1$, use exact Gaussian composition instead:
\begin{equation}\label{eq:full-calibration}
 V_{\mathrm{full}}(T;\mu)=T/\mu^2.
\end{equation}
The raw expression in \eqref{eq:clt-calibration} evaluated at $q=1$ is
larger than \eqref{eq:full-calibration} and is not the exact Gaussian formula.
Write $V_{\mathrm{model}}$ for $V_{\mathrm{CLT}}$ on $(0,1)$ and
$V_{\mathrm{full}}$ at one. With this positive variance, the following
convergence theorem applies exactly; for $q<1$, an additional privacy
argument is required to assert finite $(\varepsilon,\delta)$-DP.

\begin{theorem}[Expected-gradient convergence]\label{thm:main}
Consider \eqref{eq:algorithm} at any fixed $V>0$, in particular
$V=V_{\mathrm{model}}$.
Fix a dataset and suppose that $f:\R^d\to\R$
is differentiable, $\nabla f(x)=N^{-1}\sum_i h_i(x)$, and, for every $x,y$,
\begin{equation}\label{eq:assumptions}
 \norm{\nabla f(x)-\nabla f(y)}\le L\norm{x-y},\qquad
 \frac1N\sum_i\norm{h_i(x)-\nabla f(x)}^2\le\sigma_{\sg}^2.
\end{equation}
Assume $L>0$, $\sigma_{\sg}\ge0$, $f\ge f_{\mathrm{lower}}$, and let
$F\ge F_{\gap}:=f(X_0)-f_{\mathrm{lower}}$. If
\begin{equation}\label{eq:step-domain}
 C\ge4\sigma_{\sg},\qquad
 b=\frac{1-q}{B},\qquad 0<\eta\le\bar\eta:=\frac1{L(1+b)},
\end{equation}
then
\begin{equation}\label{eq:utility}
 \frac1T\sum_{t=0}^{T-1}\E\norm{\nabla f(X_t)}
 \le U(\eta):=\sqrt{D(\eta)}+\frac{D(\eta)}C,
 \qquad
 D(\eta)=\frac{4F}{\eta T}+\frac{8\sigma_{\sg}^4}{C^2}
            +2\eta L\left(b\sigma_{\sg}^2+\frac{C^2dV}{B^2}\right).
\end{equation}
The expectation is over the complete sampling and Gaussian history.
\end{theorem}

\paragraph{Proof idea.}
The Poisson second moment contains an extra squared-mean term from the random
batch count. The stepsize condition absorbs that term. A clipping-alignment
inequality then gives a telescoping bound on
$\min\{\norm{\nabla f}^2,C\norm{\nabla f}\}$, which implies
\eqref{eq:utility}. Appendix~\ref{sec:utility-proof} gives the proof.
Privacy calibration enters this argument only through the chosen variance.

\begin{corollary}[Learning rate minimizing the bound]\label{cor:rate}
Fix $q,T,C,L,\sigma_{\sg},F,V$ as above, with $F>0$, and define
\[
 W=b\sigma_{\sg}^2+\frac{C^2dV}{B^2}>0.
\]
For any fixed public cap $0<\eta_{\max}\le\bar\eta$, the choice
\begin{equation}\label{eq:eta-star}
 \eta_\star=\min\left\{\eta_{\max},
            \sqrt{\frac{2F}{TLW}}\right\}
\end{equation}
satisfies $U(\eta_\star)\le U(\eta)$ for every $0<\eta\le\eta_{\max}$.
This holds in particular at $V=V_{\mathrm{model}}$; the minimization is
exact for that noise choice, independently of its privacy interpretation.
\end{corollary}

\begin{proof}
The map $z\mapsto\sqrt z+z/C$ is strictly increasing on $[0,\infty)$,
so it suffices to minimize $D$. Set $A=4F/T$, $J=2LW$, and
$D_0=8\sigma_{\sg}^4/C^2$. Then
\[
 D(\eta)=A/\eta+D_0+J\eta,\qquad
 D'(\eta)=-A/\eta^2+J.
\]
It decreases until $\eta_0=\sqrt{A/J}$ and increases afterwards.
The minimum on $(0,\eta_{\max}]$ is therefore attained at
$\min\{\eta_{\max},\eta_0\}$, which is \eqref{eq:eta-star}.
\end{proof}

Thus, when the cap is inactive,
\begin{equation}\label{eq:optimal-value}
 D(\eta_\star)=\frac{8\sigma_{\sg}^4}{C^2}
                     +4\sqrt{\frac{2FLW}{T}}.
\end{equation}
When the cap is active, substitute $\eta_{\max}$ into \eqref{eq:utility}.
A more conservative cap, such as $1/[2L(1+b)]$, changes only the first
argument of the minimum in \eqref{eq:eta-star}.

\paragraph{Public selection and the zero-gap case.}
The optimizer of an upper bound need not be the best learning rate for actual
training performance. Also, substituting a privately computed exact objective
gap into \eqref{eq:eta-star} would require accounting for that private selection.
To use the theorem as a private parameter-selection rule, $F,L$ and
$\sigma_{\sg}$ must be public inputs common across neighbors, with the
utility assumptions valid for the dataset under consideration. Public uniform
bounds on a dataset class are one sufficient choice. The theorem does not
require those bounds to be attained.

If $F=0$, then $D(\eta)=D_0+2\eta LW$, where $D_0=8\sigma_{\sg}^4/C^2$.
Both $D$ and $U$ are strictly increasing for $\eta>0$, so neither has a
minimizing positive learning rate. On any admissible interval
$(0,\eta_{\max}]$,
\[
 \inf D(\eta)=D_0,\qquad
 \inf U(\eta)=\sqrt{D_0}+D_0/C.
\]
Both infima are approached as $\eta\downarrow0$. Positive chosen
variance and $d\ge1$ ensure $W>0$.

\section{Batch-size predictions of the analytical model}
\label{sec:clt-batches}
Fix the public privacy target, hence $\mu$, and the problem parameters
$N,d,C,L,\sigma_{\sg},F$, with $F>0$. Substitution in \eqref{eq:utility}
gives, for $0<q<1$,
\begin{equation}\label{eq:clt-D}
 D_{\mathrm{CLT}}(\eta)
 =\frac{4F}{\eta T}+D_0
  +2\eta L\left[\frac{1-q}{qN}\sigma_{\sg}^2+
  \frac{C^2d}{N^2q^2\log(1+\mu^2/(Tq^2))}\right],
 \qquad D_0=\frac{8\sigma_{\sg}^4}{C^2}.
\end{equation}
The optimal rate remains \eqref{eq:eta-star}, with the bracketed quantity
as $W$. In particular, changing the privacy prescription changes both the
noise term and its minimizing learning rate.

\paragraph{The quantity that survives learning-rate minimization.}
When the cap is inactive, \eqref{eq:optimal-value} depends on $W/T$.
Its privacy contribution is $C^2dR/N^2$, where
\begin{equation}\label{eq:clt-normalization}
 R_{\mathrm{CLT}}(T,q)
 =\frac{V_{\mathrm{CLT}}}{Tq^2}
 =\frac{1}{m\log(1+\mu^2/m)},\qquad m=Tq^2.
\end{equation}
If the sampling contribution $S=b\sigma_{\sg}^2$ to $W$ is negligible, the minimized coefficient above $D_0$
is approximately
\[
 \frac{4C}{N}\sqrt{2FLd\,R_{\mathrm{CLT}}(T,q)}.
\]
This is an approximation to the bound, not an upper bound obtained by
discarding a positive term. Proposition~\ref{prop:dominance} makes it
quantitative: if $S$ is at most $\rho$ times the Gaussian-noise contribution $P$,
the two minimized coefficients above $D_0$ differ by at most a factor
$\sqrt{1+\rho}$, including active caps. That proposition is independent of
the choice of privacy calibration.

\begin{lemma}[Batch comparison for a prescribed variance]\label{lem:batch-general}
Fix the problem parameters with $F>0$, prescribe $V(T,q)>0$, and use the
full admissible cap.
At fixed $T$, if $V(T,q)/q^2$ is nonincreasing in $q$, then the minimized
bound is nonincreasing, including active caps.
At exact expected passes $Tq=E$, if $V(T,q)/(Tq^2)$ is nonincreasing
as $q$ increases along the compared runs, then the minimized bound is nonincreasing
whenever both unconstrained rates are admissible.
\end{lemma}
\begin{proof}
At fixed $T$, $W$ is nonincreasing and $\bar\eta(q)$ is nondecreasing.
Every previously admissible rate remains admissible with no larger
$D(\eta)$, so minimization preserves the comparison.
At exact expected passes,
\[
 \frac WT=\frac{1-q}{EN}\sigma_{\sg}^2+
 \frac{C^2d}{N^2}\frac{V(T,q)}{Tq^2}.
\]
This is nonincreasing. With inactive caps the result follows from
$D_\star=D_0+4\sqrt{2FLW/T}$ and monotonicity of $\sqrt D+D/C$.
\end{proof}

The common monotonicity calculation below uses
$h_a(s)=s\log(1+a/s)$, which is increasing in both $a,s>0$:
\begin{equation}\label{eq:log-perspective}
 \partial_a h_a(s)=\frac{s}{s+a}>0,\qquad
 h_a'(s)=\log(1+a/s)-\frac{a}{s+a}>0.
\end{equation}
The second inequality follows by integrating the positive derivative
$z/(1+z)^2$ of $\log(1+z)-z/(1+z)$ from zero to $z=a/s$.

\begin{proposition}[Batch comparisons within the model]\label{prop:clt-batch}
Use $V=V_{\mathrm{model}}$ and the full cap $\bar\eta(q)$ from
\eqref{eq:step-domain}. At fixed integer $T\ge1$, the minimized upper bound
is nonincreasing in $q\in(0,1]$.
If instead integer horizons satisfy $Tq=T'q'=E>0$ and $q\le q'$, then
$W(T',q')/T'\le W(T,q)/T$. The minimized upper bound is therefore
nonincreasing when both unconstrained rates satisfy their caps.
These are comparisons of finite convergence bounds at the chosen noise;
they do not supply finite privacy for the subsampled runs.
\end{proposition}
\begin{proof}
For fixed $T$, take $a=\mu^2/T$ and $s=q^2$ in
\eqref{eq:log-perspective}. Thus $V_{\mathrm{CLT}}/q^2=1/h_a(q^2)$
decreases, and Lemma~\ref{lem:batch-general} gives the comparison.

At exact expected passes, $m=Eq$ and
\begin{equation}\label{eq:clt-fixed-passes}
 \frac{W(T,q)}T
 =\frac{1-q}{EN}\sigma_{\sg}^2+
   \frac{C^2d}{N^2}\frac{1}{Eq\log(1+\mu^2/(Eq))}.
\end{equation}
Both terms are nonincreasing by \eqref{eq:log-perspective}, with
$a=\mu^2$ and $s=Eq$. Lemma~\ref{lem:batch-general} applies. At $q=1$, exact Gaussian
composition gives $R=1/\mu^2$, which is smaller than every raw CLT value
because $\log(1+x)<x$. This also gives the endpoint comparison in both
budget policies.
\end{proof}

\paragraph{Why an interior best batch can still occur.}
At fixed expected passes, write $\tau=T\eta$. The coefficient becomes
$4F/\tau+D_0+2\tau L(W/T)$, but its cap is
\[
 \tau\le\frac{EN}{L[1+(N-1)q]}.
\]
For $N>1$, this admissible interval shrinks with $q$. Thus improving the
privacy contribution does not imply an unconditional ordering of the
optimized bounds. With $T=\lceil E/q\rceil$, one must additionally retain
the rounding slack $E\le Tq<E+q$. Our numerical comparison uses these
integer horizons and the cap, without replacing them by $T=E/q$.

\paragraph{What changes relative to Abadi scaling.}
Theorem~1 of Abadi et al.~\cite{abadi} supplies constants $c_1,c_2$
and, within $\varepsilon<c_1Tq^2$, the sufficient prescription
$V_A=c_2^2Tq^2\log(1/\delta)/\varepsilon^2$. This makes
$V_A/(Tq^2)=c_2^2\log(1/\delta)/\varepsilon^2$ independent of batch size.
By contrast, \eqref{eq:clt-normalization} varies with $Tq^2$. For large
$m=Tq^2$ it approaches $1/\mu^2$, but as $m\downarrow0$ it grows like
\[
 R_{\mathrm{CLT}}\sim\frac{1}{m\log(\mu^2/m)}.
\]
These are asymptotics of the analytical formula, not a privacy theorem in
that limit. They exhibit a small-batch penalty in the retuned analytical-model
bound that disappears under the linear-in-$q$ prescription. In the high-noise regime,
$e^{1/V}-1\simeq1/V$, so the CLT model itself gives
$\sqrt V\simeq q\sqrt T/\mu$. Its nonlinear dependence matters when this
approximation is inadequate. The Abadi cancellation alone does not prove
that small batches are optimal under the Abadi-based full bound. Moreover,
its theorem requires $\varepsilon<c_1Tq^2$; neither that restriction nor
$c_2$ may be silently dropped or set to one. At exact expected passes,
the restriction becomes $q>\varepsilon/(c_1E)$, excluding arbitrarily
small batches at fixed $E$.

\paragraph{The asymptotic privacy assumption.}\label{par:clt-scope}
The subsampling CLT of Bu et al.~\cite[Appendix B.2, Theorem 6 and
Lemma B.1]{bu2020deepgdp} holds with fixed
positive multiplier variance and $q\sqrt T\to\nu\in(0,\infty)$. Its
Gaussian limit has $\mu=\nu\sqrt{e^{1/V}-1}$, motivating
\eqref{eq:clt-calibration}. An approximation at finite $T$ requires an
error estimate before it becomes a finite privacy guarantee. A directly
covered sequence is $q_T=\nu/\sqrt T$ and the fixed variance
$V=1/\log(1+\mu^2/\nu^2)$, for sufficiently large $T$.
In particular, at fixed $E=Tq$ and $q\downarrow0$,
\[
 q\sqrt T=\sqrt{Eq}\longrightarrow0,\qquad
 V_{\mathrm{CLT}}=\frac{1}{\log(1+\mu^2/(Eq))}\longrightarrow0.
\]
The fixed-noise CLT does not justify this moving-noise limit. In fact,
Lemma~\ref{lem:fixed-pass-clt-collapse} proves that the corresponding
scalar product experiments converge in total variation to each other.
Their limiting trade-off is $G_0=\mathrm{Id}$, rather than the prescribed
$G_\mu$ with $\mu>0$. The adaptive transport gives the same
indistinguishability conclusion for neighboring transcript laws.
Thus the fixed-pass comparisons describe the analytical noise model;
they do not establish equal actual privacy. This limiting discrepancy
is toward stronger privacy and does not prove a privacy violation in any
finite numerical cell. The fixed-$q$ limit in \emph{Subsampling is not
Magic}~\cite{snm} likewise cannot be specialized to $q=E/T$.

\section{Comparison with numerical accounting}
\label{sec:clt-numerical}
We compare the CLT model, saved numerical PLD calibration, and the finite
sufficient calibration $V_\star$ of Appendix~\ref{sec:privacy}. All three
use the same convergence bound and their own minimizing rate with the
full cap $1/[L(1+b)]$. The illustrative parameters are
\[
 N=50000,\quad d=100000,\quad C=0.4,\quad L=0.25,\quad
 \sigma_{\sg}^2=0.0099,\quad F=0.637145646205098.
\]
They satisfy $C\ge4\sigma_{\sg}$. We take $\delta=10^{-5}$,
$\varepsilon\in\{1,8\}$, and the batch grid
\[
 B\in\{64,256,512,1024,2048,4096,8192,50000\},
\]
separately at $T=1000$ and $T=\lceil8/q\rceil$.
At full sampling the CLT-model curve uses
\eqref{eq:full-calibration}.

The retained PLD values use Google's \texttt{dp-accounting} version 0.6.0
for Poisson subsampling and add/remove adjacency, with pessimistic
discretization, connect-the-dots, mesh
$\min(10^{-3},0.1\varepsilon/T)$, log-mass truncation $-50$, composition
tail truncation $10^{-15}$, and inversion tolerance $10^{-7}$.
These are numerical reference values, not exact minimum variances or
interval-certified computations. A fresh run of the same accountant reevaluates
all 32 saved multipliers and finds $\delta_{\rm upper}\le10^{-5}$ in every
case. Three cells are also reinverted at the original mesh and half that mesh,
as reported below. The curves retain the original calibrations; the refinement
is a separate sensitivity check. No training outcomes are measured.

\begin{table}[ht]
\centering\small
\input{dpsgd_analytic_gaussian_bound.evidence/clt-model/comparison-table.tex}
\caption{Variance ratios over the subsampled grid and batch sizes minimizing
the displayed bound on the eight-point grid. ``Finite'' means $V_\star$.
Each method uses its own capped optimal learning rate.}
\label{tab:clt-comparison}
\end{table}

The CLT prescription selects less noise than the saved PLD calibration in
every subsampled cell. The largest difference occurs at $T=1000$,
$\varepsilon=1$, where the variance ratio is approximately $0.589$.
This variance comparison does not itself prove a privacy violation: the PLD values
are numerical sufficient calibrations, not lower bounds on the exact
required variance. It does show that replacing them by the CLT prescription
is a substantive modeling choice. For $T=\lceil8/q\rceil$ and
$\varepsilon=1$, the best grid batch changes from $8192$ under PLD to
$4096$ under the CLT model. The respective bounds must not be interpreted
as already establishing equal finite privacy. For the most discrepant cell,
$T=1000$, $B=64$, $\varepsilon=1$, Proposition~\ref{prop:finite-clt-failure}
gives a separate event-based proof of failure of the target transcript privacy
for a constant query. This is a lower-bound argument, independent of the
pessimistic accountant's upper estimates.

\paragraph{Numerical sensitivity and winner margins.}
Reinversion at the original mesh reproduces the three selected saved
multipliers within the $10^{-7}$ inversion tolerance. At half the mesh,
the multipliers for $(T,B)=(1000,64),(98,4096),(49,8192)$ at
$\varepsilon=1$ are approximately $0.667837898$, $3.254955024$ and
$4.526649890$. The accountant's upper estimates of $\delta$ at the CLT
noise are respectively $2.15208\times10^{-3}$, $2.17605\times10^{-5}$
and $2.00916\times10^{-5}$. Exceeding the target in this sufficient test
does not alone prove privacy failure.

For the two rounded-pass candidates $B=4096,8192$, the refined PLD bounds
are $0.2363835891$ and $0.2355537135$, preserving the original ordering.
The relative excess over the smaller value is about $0.35\%$; the CLT
ordering of the same two batches has a margin of about $0.23\%$.
For the finite sufficient calibration, the runner-up bound is only about
$0.00010\%$ above the winning bound.
These are small differences between grid values of an upper bound, not
evidence of a substantial difference in training performance. The retained
script, settings, input hashes and fresh outputs make these checks reproducible.

The sampling term $S=b\sigma_{\sg}^2$ retained in $W$ is at most
$0.00481$ times the Gaussian-noise term $P=C^2dV/B^2$ across the
CLT-model grid. The complete conditional Poisson sampling variance is
$b(\norm m^2+\nu^2)$; the additional squared-mean contribution is
absorbed through the stepsize restriction. Nevertheless the stepsize cap changes the batch-size
prediction: at eight expected epochs it is active for CLT batches $8192$
and $50000$ when $\varepsilon=1$, and also for $4096$ when
$\varepsilon=8$. Privacy domination therefore does not justify ignoring
the stepsize restriction.

\begin{figure}[ht]
\centering
\includegraphics[width=0.96\textwidth]{dpsgd_analytic_gaussian_bound.evidence/clt-model/bound-comparison.pdf}
\caption{The minimized finite convergence bound at three noise choices.
The CLT curve has an approximate privacy interpretation for $q<1$;
the finite sufficient curve has the separate guarantee of
Theorem~\ref{thm:privacy}. Fixed updates and rounded expected epochs give
different comparisons.}
\label{fig:clt-comparison}
\end{figure}

\clearpage
\appendix
\section{A separate finite Gaussian-reference guarantee}\label{sec:privacy}

The calibration here is sufficient for finite privacy. It is distinct from
$V_{\mathrm{model}}$ and is not claimed to be equally accurate as an
approximation to numerical accounting. The same convergence theorem and
learning-rate optimizer apply after substituting this variance.

\paragraph{Explicit privacy calibration.}
For $\varepsilon>0$ and $0<\delta<1$, define
\begin{equation}\label{eq:order}
 \ell=\log(1/\delta),\quad M=Tq^2,\quad
 \theta=\varepsilon/2,\quad a=1-e^{-\theta},\quad
 p=\left\lceil\frac{\varepsilon+\ell+\log2}{\log2}\right\rceil.
\end{equation}
Since $\varepsilon+\ell>0$, the integer $p$ is at least two. For $q<1$ put
\begin{equation}\label{eq:residual-variance}
 \lambda_0=\log(1+a^2/4),\qquad
 r_0=\sqrt{\frac{2\lambda_0}{M(1-q)^2}},\qquad
 V_R=\frac{p-1}{\log(1+r_0)}.
\end{equation}
Set $V_R=0$ at $q=1$. For all $q\in(0,1]$, put
\begin{equation}\label{eq:gaussian-variance}
 m_0=\frac{\varepsilon}{\sqrt{2\varepsilon+2\ell}+\sqrt{\varepsilon+2\ell}},
 \qquad V_G=\frac{M}{m_0^2},\qquad V_\star=\max\{V_R,V_G\}.
\end{equation}
All logarithms are natural. The subscripts $R$ and $G$ distinguish the
residual and Gaussian requirements.

\begin{theorem}[Finite privacy]\label{thm:privacy}
Use $V\ge V_\star$ in \eqref{eq:algorithm}, with all algorithm parameters
fixed publicly across neighboring datasets. Then its transcript and final
iterate satisfy $(\varepsilon,\delta)$-differential privacy.
\end{theorem}



\subsection{A common scalar experiment for adaptive releases}
For $V>0$, let
\[
 Q_V=\N(0,V),\qquad P_V=(1-q)\N(0,V)+q\N(1,V).
\]
For probability laws $P,Q$, write
$H_{e^u}(P\mid Q)=\sup_A\{P(A)-e^uQ(A)\}$, where $A$ ranges over
measurable sets. Its density expression is $\int(p-e^u q)_+$.
Define the two-direction product profile
\[
 \Delta(V)=\max\{H_{e^\varepsilon}(P_V^{\otimes T}\mid Q_V^{\otimes T}),
                 H_{e^\varepsilon}(Q_V^{\otimes T}\mid P_V^{\otimes T})\}.
\]

\begin{lemma}[Adaptive transport]\label{lem:transport}
Both directed privacy errors of the transcript in \eqref{eq:algorithm}
are at most $\Delta(V)$.
\end{lemma}
\begin{proof}
Fix neighboring datasets. At any common transcript history their public
recursion reconstructs the same state. Write the differing clipped query as
$v=Cw$, where $\norm w\le1$. Given a scalar input $z$, generate
\[
 Cwz+\xi,\qquad \xi\sim\N(0,C^2V(I_d-ww^{\mathsf T})).
\]
The covariance is positive semidefinite. Under $Q_V$ the output is
$\N(0,C^2VI_d)$; under $P_V$ it is that Gaussian plus an independent
Bernoulli-$q$ multiple of $v$. Simulate the common slots independently,
add their contributions, and advance the public recursion. Iterating gives
one measurable randomized map from the scalar product experiment to the
two transcript laws. It is the same map on both sides of this fixed
neighboring pair. Data processing for $H_{e^\varepsilon}$ proves the claim.
This argument does not require independence of the optimizer's privacy-loss
increments.
\end{proof}

\subsection{The Gaussian reference and its residual}
Standardize by $u=V^{-1/2}$ and set
\[
 Q=\N(0,1),\qquad P=(1-q)\N(0,1)+q\N(u,1),\qquad G=\N(qu,1).
\]
Write $P_T=P^{\otimes T}$, $Q_T=Q^{\otimes T}$ and $G_T=G^{\otimes T}$.
The two directions between $G_T$ and $Q_T$ have the exact
Gaussian profile $g_{\mu_{\mathrm{ref}}}$ from \eqref{eq:gaussian-profile}, with the
reference separation
\[
 \mu_{\mathrm{ref}}=q\sqrt{T/V}.
\]
This reference separation differs from the CLT parameter in
\eqref{eq:clt-calibration}.
Indeed, the product Gaussian likelihood depends only on the coordinate sum;
standardizing that sum gives separation $\mu_{\mathrm{ref}}$. Thresholding the likelihood
ratio and integrating gives \eqref{eq:gaussian-profile}; reflection exchanges
the two directions. The profile is increasing in $\mu_{\mathrm{ref}}$, by adding Gaussian
noise and rescaling.

\begin{lemma}[Centered moment bound]\label{lem:moment}
For every integer $p\ge2$, the density ratio
$R=dP_T/dG_T$ satisfies
\begin{equation}\label{eq:moment}
 \E_{G_T}|R-1|^p\le H^p,\qquad
 H=(e^S-1)^{1/2},\quad
 S=\frac{Tq^2(1-q)^2}{2}x^2e^x,\quad x=\frac{p-1}{V}.
\end{equation}
\end{lemma}
\begin{proof}
For $A\sim\operatorname{Bernoulli}(q)$, the one-step ratio in coordinates
centered at $G$ is
\[
 r_u(z)=\E_A\exp\{u(A-q)z-u^2(A-q)^2/2\},\qquad
 R_u(z_1,\ldots,z_T)=\prod_{t=1}^T r_u(z_t).
\]
Let $\gamma_k$ be standard Gaussian measure on $\R^k$ for $k\ge1$, and define
\[
 (\mathcal U_\rho F)(z)=\E_Z F(\rho z+\sqrt{1-\rho^2}Z),
 \qquad Z\sim\gamma_T.
\]
Gaussian hypercontractivity~\cite[p.~1062]{gross} states that
$\|\mathcal U_\rho F\|_{L^p(\gamma_T)}\le\|F\|_{L^2(\gamma_T)}$
for $F\in L^2(\gamma_T)$, $2\le p<\infty$, and $0\le\rho\le(p-1)^{-1/2}$.
Here $R_u-1$ is a finite signed sum of Gaussian shift likelihoods,
so all its finite Gaussian moments exist.
Gaussian integration sends each shift vector $v$ to $\rho v$, since
\[
 \mathcal U_\rho e^{\langle v,\cdot\rangle-\|v\|^2/2}(z)
 =e^{\langle\rho v,z\rangle-\|\rho v\|^2/2}.
\]
All shifts in $R_u$ are linear in $u$. Thus, for
$w=\sqrt{p-1}\,u$ and $\rho=(p-1)^{-1/2}$,
\begin{equation}\label{eq:hypercontractivity}
 R_u-1=\mathcal U_\rho(R_w-1),\qquad
 \E|R_u-1|^p\le\bigl(\E(R_w-1)^2\bigr)^{p/2}.
\end{equation}
For independent $A,A'\sim\operatorname{Bernoulli}(q)$, put
\[
 K=\E e^{w^2(A-q)(A'-q)}.
\]
Gaussian integration and independence give
\[
 \E_{\gamma_1}r_w=1,\qquad \E_{\gamma_1}r_w^2=K,\qquad
 \E_{\gamma_T}(R_w-1)^2=K^T-1.
\]
The variable $Z_0=w^2(A-q)(A'-q)$ satisfies
\[
 \E Z_0=0,\qquad Z_0\le w^2,\qquad
 \E Z_0^2=w^4q^2(1-q)^2.
\]
Taylor's integral remainder therefore yields
\[
 K-1\le\tfrac12q^2(1-q)^2w^4e^{w^2},\qquad
 K^T-1\le\exp\!\left(\tfrac12Tq^2(1-q)^2w^4e^{w^2}\right)-1.
\]
Combining this with \eqref{eq:hypercontractivity} and $w^2=(p-1)/V$
proves the claim.
\end{proof}

\begin{lemma}[Both directed privacy errors]\label{lem:directions}
Put $\theta=\varepsilon/2$ and $c_p=(p-1)^{p-1}/p^p$. Then
\begin{equation}\label{eq:budget}
 \Delta(V)\le e^\theta g_{\mu_{\mathrm{ref}}}(\theta)
       +\frac{e^\varepsilon c_pH^p}{(1-e^{-\theta})^{p-1}}.
\end{equation}
\end{lemma}
\begin{proof}
Write $E_p=c_pH^p/(e^\theta-1)^{p-1}$.
For $a_0>0$, maximizing $(z-a_0)_+/z^p$ over $z>0$ gives
\[
 (z-a_0)_+\le\frac{c_p(z_+)^p}{a_0^{p-1}},\qquad z\in\R.
\]
Apply this to $R-1$ and $1-R$, with thresholds $e^\theta-1$ and
$1-e^{-\theta}$, respectively, using $(\pm(R-1))_+^p\le|R-1|^p$.
Lemma~\ref{lem:moment} implies, for every measurable event $A$,
\[
 P_T(A)\le e^\theta G_T(A)+E_p,\qquad
 G_T(A)\le e^\theta P_T(A)+e^{p\theta}E_p.
\]
For the reverse inequality use
$(1-e^\theta R)_+=e^\theta(1-R-(1-e^{-\theta}))_+$.
Compose the forward inequality with the Gaussian bound from $G_T$ to
$Q_T$ at parameter $\varepsilon-\theta$. In the other direction first
compare $Q_T$ to $G_T$, and then $G_T$ to $P_T$. The latter residual
therefore receives the factor $e^{\varepsilon-\theta}e^{p\theta}$.
The two directed errors are consequently bounded by
\[
 e^\theta g_{\mu_{\mathrm{ref}}}(\theta)+E_p,\qquad
 g_{\mu_{\mathrm{ref}}}(\theta)+e^{\varepsilon+(p-1)\theta}E_p.
\]
Both are at most the right-hand side of \eqref{eq:budget}: $g_{\mu_{\mathrm{ref}}}(\theta)\ge0$, and the larger
residual equals $e^\varepsilon c_pH^p/(1-e^{-\theta})^{p-1}$.
\end{proof}

\subsection{Verification of the calibration}
Use \eqref{eq:order}--\eqref{eq:gaussian-variance}. For $r\ge0$,
\begin{equation}\label{eq:log-inequality}
 \log(1+r)\sqrt{1+r}\le r.
\end{equation}
To see this, put $w=\sqrt{1+r}\ge1$. The function
$w-w^{-1}-2\log w$ vanishes at one and has derivative $(w-1)^2/w^2$.
When $q<1$ and $V\ge V_R$, we have $0\le x\le\log(1+r_0)$, hence
$x^2e^x\le r_0^2$. It follows that $S\le\lambda_0$ and $H\le a/2$.
The residual in \eqref{eq:budget} obeys
\begin{align*}
 \frac{e^\varepsilon c_pH^p}{a^{p-1}}
 &\le e^\varepsilon a\,2^{-p}
 \le e^\varepsilon2^{-p}\le\delta/2.
\end{align*}
Here $c_p\le1$, $0<a<1$, and the last inequality is exactly the order
choice in \eqref{eq:order}. At $q=1$, $P=G$ and the residual is zero.

For the Gaussian part, put $t=\sqrt{\varepsilon+2\ell}$. The definition
of $m_0$ gives $\varepsilon/(2m_0)-m_0/2=t$.
If $V\ge V_G$, then $\mu_{\mathrm{ref}}\le m_0$. Monotonicity of the Gaussian profile,
followed by dropping its nonnegative subtracted term, yields
\[
 e^\theta g_{\mu_{\mathrm{ref}}}(\varepsilon-\theta)
 \le e^{\varepsilon/2}\Phi(-t)
 \le\tfrac12 e^{\varepsilon/2-t^2/2}=\delta/2.
\]
The tail inequality follows by writing
$\Phi(-t)=\int_0^\infty\varphi(t+z)\,dz
\le e^{-t^2/2}\int_0^\infty\varphi(z)\,dz$ for $t\ge0$.
Both terms in \eqref{eq:budget} are now at most $\delta/2$.
Lemma~\ref{lem:transport} proves Theorem~\ref{thm:privacy}.
The formula deliberately retains the sufficient value $V_G$ at full
sampling; it does not claim the smallest feasible Gaussian variance.

\section{Convergence for any positive noise variance}\label{sec:utility-proof}
The following proof holds at any fixed noise variance $V>0$.
Thus it applies to the recursion at any $V>0$, including $V=V_\star$.

\subsection{Exact sampling moments and clipping alignment}
Let $\mathcal F_t$ be the history before update $t$, and write
$\E_t[\cdot]=\E[\cdot\mid\mathcal F_t]$.
Condition on this history, so $X_t=x$ is fixed for the fresh update. Abbreviate
\begin{align*}
 g&=\nabla f(x),\qquad r=\norm g,\qquad v_i=\clip_C(h_i(x)),\\
 m&=N^{-1}\sum_i v_i,\qquad \nu^2=N^{-1}\sum_i\norm{v_i-m}^2,
 \qquad \widehat g=B^{-1}\sum_i B_{t,i}v_i.
\end{align*}
Independence and centering of the Bernoulli indicators give
\begin{equation}\label{eq:poisson}
 \E_t\widehat g=m,\qquad
 \E_t\norm{\widehat g-m}^2=b\bigl(\norm m^2+\nu^2\bigr),\qquad
 \E_t\norm{\widehat g}^2=(1+b)\norm m^2+b\nu^2.
\end{equation}
The squared-mean term remains even when all $v_i$ are identical.
It is therefore essential to the stepsize restriction.

By variance decomposition and nonexpansiveness of projection,
\begin{equation}\label{eq:clip-moments}
 \nu^2+\norm{m-\clip_C(g)}^2
 =\frac1N\sum_i\norm{v_i-\clip_C(g)}^2
 \le\frac1N\sum_i\norm{h_i(x)-g}^2\le\sigma_{\sg}^2.
\end{equation}

\begin{lemma}[Clipping alignment]\label{lem:alignment}
If $C\ge4\sigma_{\sg}$ and $\phi_C(r)=\min\{r^2,Cr\}$, then
\begin{equation}\label{eq:alignment}
 \langle g,m\rangle\ge\tfrac14\phi_C(r)+\tfrac12\norm m^2
                          -\frac{2\sigma_{\sg}^4}{C^2}.
\end{equation}
\end{lemma}
\begin{proof}
If $r<C/2$, clipping a record requires
$\norm{h_i(x)-g}>C-r$. Its clipping displacement is at most
$\norm{h_i(x)-g}$, because $g$ belongs to the clipping ball. Since
$N^{-1}\sum_i(h_i(x)-g)=0$, we obtain
\[
 \norm{m-g}\le N^{-1}\sum_i\norm{h_i(x)-g}
  \mathbf1_{\{\norm{h_i(x)-g}>C-r\}}
 \le\frac{\sigma_{\sg}^2}{C-r}\le\frac{2\sigma_{\sg}^2}{C}.
\]
Use $2\langle g,m\rangle=r^2+\norm m^2-\norm{m-g}^2$.
If $C/2\le r\le C$, \eqref{eq:clip-moments} gives
$\norm{m-g}^2\le\sigma_{\sg}^2\le r^2/4$; the same identity gives
$\langle g,m\rangle\ge3r^2/8+\norm m^2/2$.
Finally, if $r>C$, apply the identity to $v=Cg/r$ and $m$:
\[
 \langle v,m\rangle\ge\tfrac12(C^2+\norm m^2-\sigma_{\sg}^2)
 \ge\tfrac{15}{32}C^2+\tfrac12\norm m^2.
\]
Multiplying by $r/C>1$ proves \eqref{eq:alignment} in the last regime.
All three estimates imply the stated conservative common bound.
\end{proof}

\subsection{Descent, telescoping, and the expected gradient norm}
Let $z_t=C\sqrt V Z_t/B$. Global smoothness, zero mean of $z_t$, and its
independence from the sample give
\begin{align}
 \E_t f(X_{t+1})
 &\le f(x)-\eta\langle g,m\rangle
   +\frac{L\eta^2}{2}\left((1+b)\norm m^2+b\nu^2+s^2\right)\notag\\
 &\le f(x)-\frac\eta4\phi_C(r)+\frac{2\eta\sigma_{\sg}^4}{C^2}
       +\frac{L\eta^2}{2}(b\sigma_{\sg}^2+s^2)\notag\\
 &\hspace{2em}-\frac\eta2[1-\eta L(1+b)]\norm m^2.\label{eq:descent}
\end{align}
The final term is nonpositive by \eqref{eq:step-domain}.
All expectations are finite: $\norm{\widehat g}\le C/q$ and Gaussian
noise has finite second moment, so each finite-time iterate has finite second
moment. Global smoothness gives linear gradient growth and at most quadratic
growth of $|f|$. In particular, the conditional estimates and the following
telescoping are legitimate.

Take expectations in \eqref{eq:descent}, sum over $0\le t<T$, and use
$f(X_T)\ge f_{\mathrm{lower}}$ and $F_{\gap}\le F$. This gives
\begin{equation}\label{eq:progress}
 \frac1T\sum_{t<T}\E\phi_C(\norm{\nabla f(X_t)})\le D(\eta).
\end{equation}
For every $r\ge0$,
$r\le\sqrt{\phi_C(r)}+\phi_C(r)/C$.
Applying Jensen's inequality jointly over the uniformly weighted times and
the algorithm's randomness proves \eqref{eq:utility}.
This completes Theorem~\ref{thm:main}. Its constants are sufficient rather
than optimal. It bounds pre-update expected gradient norms; it makes no
corresponding final-iterate gradient claim.

\section{Batch comparisons for the finite calibration}
\label{sec:batch}
Keep $N,d,C,L,\sigma_{\sg},F,\varepsilon,\delta$ fixed, with $F>0$,
and choose the public variance $V=V_\star(T,q)$ from
\eqref{eq:gaussian-variance}. Write
\[
 P(q)=\frac{C^2dV_\star(T,q)}{N^2q^2},\qquad
 S(q)=\frac{1-q}{qN}\sigma_{\sg}^2,\qquad W(q)=S(q)+P(q).
\]
Use the full admissible cap $\bar\eta(q)=1/[L(1+(1-q)/(qN))]$,
and let $\eta_\star(q)$ be \eqref{eq:eta-star} with this cap and $W(q)$.
Denote the resulting upper bound by $U_\star(q)=U(\eta_\star(q))$.

\begin{theorem}[Fixed-update batch comparison]\label{thm:batch}
For fixed integer $T\ge1$ and $0<q\le q'\le1$,
\[
 \frac{V_\star(T,q')}{q'^2}\le\frac{V_\star(T,q)}{q^2},
 \qquad U_\star(q')\le U_\star(q).
\]
Under the assumptions of Theorem~\ref{thm:main}, the publicly tuned algorithm
at $q'$ is $(\varepsilon,\delta)$-DP and its average expected gradient norm
is at most $U_\star(q')\le U_\star(q)$. The comparison includes an active
stepsize cap.
\end{theorem}
\begin{proof}
The Gaussian contribution satisfies $V_G/q^2=T/m_0^2$, independently of $q$.
For the residual contribution, put $k=\sqrt{2\lambda_0/T}>0$. For $0<q<1$,
\[
 \frac{V_R(T,q)}{q^2}
 =\frac{p-1}{q^2\log(1+k/[q(1-q)])}.
\]
By \eqref{eq:log-perspective}, $h_a(s)$ increases in both arguments.
Set $s=q^2$ and $a=kq/(1-q)$. Both increase with $q$, and
\[
 s\log(1+a/s)=q^2\log(1+k/[q(1-q)]).
\]
Thus $V_R/q^2$ is nonincreasing on $(0,1)$. At $q=1$ the residual is zero,
so the same comparison holds at the endpoint. Taking the maximum with the
constant Gaussian contribution proves the first claim.

Lemma~\ref{lem:batch-general} now proves the optimized-bound comparison.
The privacy and expected-gradient claims follow from
Theorems~\ref{thm:privacy} and~\ref{thm:main} for that same public run.
\end{proof}

This theorem compares upper bounds at equal update counts. Larger batches
then use more expected record evaluations. It neither orders actual training
errors nor supplies an equal-computation comparison.

\begin{proposition}[Quantitative privacy domination]\label{prop:dominance}
Fix a batch probability and a public cap $0<\eta_{\max}\le\bar\eta$.
Let $P=C^2dV/B^2>0$ and $S=b\sigma_{\sg}^2\ge0$ be the
Gaussian-noise and sampling contributions to $W=P+S$, and let
$D_0=8\sigma_{\sg}^4/C^2$. Define the two minimized coefficients
\[
 D_{\mathrm{priv}}=\min_{0<\eta\le\eta_{\max}}
     \left\{\frac{4F}{\eta T}+D_0+2\eta LP\right\},\qquad
 D_{\mathrm{full}}=\min_{0<\eta\le\eta_{\max}}
     \left\{\frac{4F}{\eta T}+D_0+2\eta L(P+S)\right\}.
\]
If $S\le\rho P$ for $\rho\ge0$, then, including an active cap,
\[
 D_{\mathrm{priv}}-D_0
 \le D_{\mathrm{full}}-D_0
 \le\sqrt{1+\rho}\,(D_{\mathrm{priv}}-D_0).
\]
\end{proposition}
\begin{proof}
The lower inequality follows by increasing the variance at each $\eta$.
For the upper inequality let $k=\sqrt{1+\rho}\ge1$, and let $e$ minimize
the privacy-only coefficient. The step $e/k$ is positive and admissible.
Since $P+S\le k^2P$, its full coefficient above $D_0$ is at most
\[
 \frac{4F}{(e/k)T}+2(e/k)Lk^2P
 =k\left(\frac{4F}{eT}+2eLP\right).
\]
Minimizing the full coefficient can only decrease this value.
\end{proof}

The privacy-only coefficient is an approximation to the convergence bound;
it is not itself a valid upper bound obtained by dropping a positive term.
The proposition quantifies this approximation for the coefficient $D-D_0$,
not directly for $U$ or the actual gradient norm.
When the privacy-only cap is inactive, its value is
\[
 D_{\mathrm{priv}}-D_0
 =\frac{4C}{N}\sqrt{2FLd\,\frac{V_\star(T,q)}{Tq^2}}.
\]
Thus the normalized quantity $V_\star/(Tq^2)$ determines this part of the
retuned bound. The residual contribution may vary with $q$; the Gaussian
contribution is the constant $1/m_0^2$.

\subsection{Fixed expected participation}
\label{sec:fixed-passes}
The effective one-step noise variance $V/q^2$ used in
\emph{Subsampling is not Magic}~\cite{snm} becomes $V/(Tq^2)$ after division
by the horizon. With exact expected participation $E=Tq$, this is
$TV/E^2$. This is also the privacy normalization in the minimized coefficient
\eqref{eq:optimal-value}. The paper's many-composition limit holds at fixed
positive $q$; it cannot be specialized to $q=E/T$ without a separate argument.
We give a finite comparison directly for the sufficient variance $V_\star$.

Now allow $T$ to vary and write
\[
 R(T,q)=\frac{V_\star(T,q)}{Tq^2},\qquad
 W(T,q)=\frac{1-q}{qN}\sigma_{\sg}^2+
             \frac{C^2dV_\star(T,q)}{N^2q^2}.
\]
At $Tq=E$, the exact identity
\begin{equation}\label{eq:fixed-pass-normalization}
 \frac{W(T,q)}{T}
 =\frac{1-q}{EN}\sigma_{\sg}^2+\frac{C^2d}{N^2}R(T,q)
\end{equation}
decomposes the coefficient retained after stepsize optimization. It is not
a formula for the full variance of an averaged gradient: even with frozen
gradients, the full sampling variance also contains the squared-mean term.
The adaptive convergence guarantee follows from Theorem~\ref{thm:main}.

\begin{theorem}[Equal expected passes]\label{thm:fixed-passes}
Let $E>0$, $T,T'\ge1$ be integers, and $0<q\le q'\le1$ satisfy
$Tq=T'q'=E$. With the other public parameters fixed as above,
\[
 R(T',q')\le R(T,q),\qquad
 \frac{W(T',q')}{T'}\le\frac{W(T,q)}{T}.
\]
If both unconstrained learning rates are admissible, that is,
\[
 \sqrt{\frac{2F}{TLW(T,q)}}\le\bar\eta(q),\qquad
 \sqrt{\frac{2F}{T'LW(T',q')}}\le\bar\eta(q'),
\]
then $U_\star(T',q')\le U_\star(T,q)$. Under the convergence assumptions,
the publicly tuned run at $(T',q')$ has finite $(\varepsilon,\delta)$ privacy
and average expected gradient norm at most this smaller upper bound.
\end{theorem}
\begin{proof}
Put $m=Tq^2$ and $\kappa=\sqrt{2\lambda_0}>0$. For $q<1$,
\[
 \frac{V_R(T,q)}{Tq^2}
 =\frac{p-1}{m\log(1+\kappa/[\sqrt m(1-q)])}.
\]
Use \eqref{eq:log-perspective} with $s=m$ and
$a=\kappa\sqrt m/(1-q)$. Its value is the displayed denominator.
Consequently the normalized residual cannot increase when both $m$ and $q$
increase. The normalized Gaussian contribution is $1/m_0^2$, and the
residual is zero at $q=1$. Taking their maximum proves the same comparison
for $R$. Since $m=Eq$, increasing $q$ increases $m$ at equal expected passes.
Lemma~\ref{lem:batch-general} proves the $W/T$ and optimized-bound comparisons.
The privacy and gradient claims use the same
publicly chosen variance and learning rate.
\end{proof}

\paragraph{Why the cap condition matters.}
For the horizon-scaled step $\tau=T\eta$, the admissible cap at equal expected
passes is
\begin{equation}\label{eq:fixed-pass-cap}
 \tau_{\max}=T\bar\eta(q)
 =\frac{EN}{L[1+(N-1)q]}.
\end{equation}
For $N>1$ this decreases as $q$ grows, whereas the bound at fixed $\tau$ is
$4F/\tau+D_0+2\tau L\,W/T$. The improving variance coefficient therefore
competes with a shrinking admissible interval.

\begin{proposition}[A capped-bound reversal]\label{prop:cap-reversal}
For every $\varepsilon>0$ and $\delta\in(0,1)$, there is a public gap
majorant $F>0$ such that, at $N=2$, $d=C=L=1$, $\sigma_{\sg}=0$, and
$V=V_\star$,
\[
 U_\star(2,1/2)<U_\star(1,1).
\]
Both runs have one expected pass. Thus the inactive-cap hypothesis cannot
simply be removed from Theorem~\ref{thm:fixed-passes}.
\end{proposition}
\begin{proof}
Let $P=W(2,1/2)>0$, $Q=W(1,1)>0$, and choose $F=1+4P+Q$.
The caps are $2/3$ and $1$, respectively, and both square-root optima exceed
their caps. Substitution gives
\[
 D_\star(2,1/2)=3F+\tfrac43P,\qquad
 D_\star(1,1)=4F+2Q.
\]
The second minus the first is $1+\tfrac83P+3Q>0$. Strict monotonicity of
$D\mapsto\sqrt D+D$ proves the claim. This orders bounds, not the realized
performance of the two algorithms.
\end{proof}

\paragraph{Integer rounding.}
The policy $T=\lceil E/q\rceil$ satisfies
$E\le Tq<E+q$; it generally does not keep expected passes exactly equal.
The proof above still compares the normalized privacy costs $R(T,q)$
whenever both $q$ and
$Tq^2$ increase, but those two conditions must be checked after rounding.
The capped-bound reversal above also follows this policy, since $1/q$ is
integer at both of its points. An inactive-cap comparison for general rounded
horizons requires a separate rounding analysis. For example, at $E=1$,
$(q,T)=(0.405,3)$ and $(q',T')=(0.5,2)$ increase both $q$ and $Tq^2$,
but $(1-q')/(T'q')=0.5>(1-q)/(Tq)\approx0.4897$.
Thus these conditions alone do not order the sampling contribution to $W/T$.


\section{The fixed-pass CLT prescription has a zero privacy limit}
\label{sec:clt-collapse}
\begin{lemma}\label{lem:fixed-pass-clt-collapse}
Fix $E,\mu>0$. For integers $T\ge\max\{1,E\}$ set
\[
 q_T=E/T,\qquad u_T=\sqrt{\log(1+\mu^2T/E^2)},\qquad
 Q=\mathcal N(0,1),\qquad
 \widetilde P_T=(1-q_T)Q+q_T\mathcal N(u_T,1).
\]
Then $\operatorname{TV}(\widetilde P_T^{\otimes T},Q^{\otimes T})\to0$.
In particular, the testing trade-off converges uniformly to
$G_0(\alpha)=1-\alpha$, not to $G_\mu$.
The same total-variation conclusion holds for neighboring adaptive
transcript laws with $q=q_T$ and $V=u_T^{-2}$, by the common transport of
Lemma~\ref{lem:transport}.
\end{lemma}
\begin{proof}
For $Z\sim Q$, put $L_T=\exp(u_T Z-u_T^2/2)$, so that
$d\widetilde P_T/dQ=1+q_T(L_T-1)$. Define
\[
 h_T^2=\E_Q\bigl(\sqrt{1+q_T(L_T-1)}-1\bigr)^2.
\]
For $x\ge-1$, rationalizing the square root gives
$(\sqrt{1+x}-1)^2\le\min\{x^2,|x|\}\le |x|^{3/2}$.
Also $|L-1|^{3/2}\le L^{3/2}+1$ for $L\ge0$.
The Gaussian exponential moment therefore gives
\[
 h_T^2\le q_T^{3/2}\bigl(\E_Q L_T^{3/2}+1\bigr)
       =q_T^{3/2}\bigl(e^{3u_T^2/8}+1\bigr).
\]
The Hellinger affinity of the one-step pair is $a_T=1-h_T^2/2\in[0,1]$.
Its product affinity is $a_T^T$, hence the product squared Hellinger
distance is $H_T^2=2(1-a_T^T)\le T h_T^2$. Consequently
\[
 H_T^2\le E^{3/2}T^{-1/2}
       \left[\left(1+\frac{\mu^2T}{E^2}\right)^{3/8}+1\right]
       =O(T^{-1/8})\longrightarrow0.
\]
For probability densities $p,r$, Cauchy--Schwarz applied to
$|p-r|=|\sqrt p-\sqrt r|(\sqrt p+\sqrt r)$ gives
$\operatorname{TV}(p,r)\le H(p,r)$ with this Hellinger convention.
Thus the product total variation tends to zero. A test with type-I
error at most $\alpha$ has type-II error at least
$1-\alpha-\operatorname{TV}$; the constant randomized test gives
$1-\alpha$ as an upper bound. This proves uniform trade-off convergence.
Finally, total variation contracts under the common Markov kernel.
\end{proof}

There is no contradiction with
$Tq_T^2(e^{u_T^2}-1)=\mu^2$. Indeed the product likelihood ratio has
second centered moment $(1+\mu^2/T)^T-1\to e^{\mu^2}-1$.
Increasingly rare likelihood-ratio tails sustain this moment even as
the experiments become indistinguishable. Matching the quadratic
quantity alone does not establish a Gaussian limit.

\section{Finite correction and calibration of CLT-model noise}
\label{sec:finite-correction}
The CLT parameter can appear in a valid finite bound, provided its correction
is retained. Here the correction is obtained by direct comparison of finite
experiments; no CLT or fixed-noise asymptotic assumption is used.

Fix an integer $T\ge1$, $0<q\le1$, $V>0$, $\varepsilon>0$, $\theta>0$, and an integer
$p\ge2$. Write
\[
 \mu_c=q\sqrt{T(e^{1/V}-1)},\qquad \mu_r=q\sqrt{T/V},\qquad
 a=\frac{Tq(1-q)}{2V},\qquad c_p=\frac{(p-1)^{p-1}}{p^p},
\]
\[
 S_p=\frac{Tq^2(1-q)^2}{2}
       \left(\frac{p-1}{V}\right)^2 e^{(p-1)/V},\qquad
 E_p=\frac{c_p(e^{S_p}-1)^{p/2}}{(e^\theta-1)^{p-1}},\qquad
 E_p^-=
 \begin{cases}0,&\theta\ge a,\\ e^{p\theta}E_p,&\theta<a.\end{cases}
\]
The complete higher-moment allowance is
\begin{equation}\label{eq:finite-correction-budget}
 \mathcal B_p(V,\theta)=\max\left\{
 e^\theta g_{\mu_r}(\varepsilon-\theta)+E_p,\quad
 g_{\mu_r}(\varepsilon-\theta)+e^{\varepsilon-\theta}E_p^-
 \right\}.
\end{equation}
For signed arguments the Gaussian profile retains its displayed formula.
One may restrict $0<\theta<\varepsilon$ in an implementation.

\begin{theorem}[Finite remainder for the adaptive transcript]
\label{thm:finite-clt-correction}
For the algorithm and adjacency in Section~\ref{sec:main}, let
$\delta_{\rm actual}(\varepsilon)$ be the supremum of
$\Pr_D(A)-e^\varepsilon\Pr_{D'}(A)$ over ordered neighboring datasets and
measurable transcript events. With the public parameters above,
\[
 \delta_{\rm actual}(\varepsilon)\le
 \min\{1,\mathcal B_p(V,\theta)\}
 \le g_{\mu_c}(\varepsilon)+\mathcal R_{\varepsilon,p,\theta}(T,q,V),
\]
where the explicit nonnegative remainder is
\begin{equation}\label{eq:finite-clt-remainder}
 \mathcal R_{\varepsilon,p,\theta}(T,q,V)
 =\left[\min\{1,\mathcal B_p(V,\theta)\}-g_{\mu_c}(\varepsilon)\right]_+.
\end{equation}
The same bound holds after public postprocessing. It covers arbitrary finite
$T,q,V$ in the stated domain, including varying-noise sequences and the full
numerical grid. The correction depends on $\varepsilon$ and the chosen public
comparison parameters; it is not claimed to be a uniform CLT error rate.
\end{theorem}
\begin{proof}
Use the scalar product pair $P_T,Q_T$ from Lemma~\ref{lem:transport} and its
mean-matched Gaussian $G_T$. Lemma~\ref{lem:moment} and the positive-part
inequality used in Lemma~\ref{lem:directions} give
\[
 P_T(A)\le e^\theta G_T(A)+E_p,\qquad
 G_T(A)\le e^\theta P_T(A)+e^{p\theta}E_p.
\]
The centered one-step likelihood ratio is
$r_u(z)=\E_A e^{u(A-q)z-u^2(A-q)^2/2}$, with
$A\sim\operatorname{Bernoulli}(q)$ and $u=V^{-1/2}$.
Jensen's inequality gives $r_u(z)\ge e^{-q(1-q)/(2V)}$.
Its product is therefore at least $e^{-a}$, so the reverse comparison has
zero additive error when $\theta\ge a$. This proves the allowance $E_p^-$.
Both Gaussian directions between $G_T$ and $Q_T$ have exact profile
$g_{\mu_r}(\varepsilon-\theta)$.
Composing the two event comparisons gives precisely
\eqref{eq:finite-correction-budget}, including both exponential factors.
The common adaptive kernel preserves both inequalities. The bound by one
is trivial; the final assertion follows from $b\le g+(b-g)_+$.
\end{proof}

\paragraph{A second-moment alternative.}
An independently available finite allowance uses the exact centered moment
\[
 A_2=(1-q)^2e^{q^2/V}+2q(1-q)e^{-q(1-q)/V}
       +q^2e^{(1-q)^2/V},\qquad M_2=A_2^T-1.
\]
Set $E_2=M_2/[4(e^\theta-1)]$, and set $E_2^-=0$ if $\theta\ge a$,
otherwise $E_2^-=e^{2\theta}E_2$. Replacing $E_p,E_p^-$ by these quantities
and $\mu_r$ by $\mu_c$ in \eqref{eq:finite-correction-budget} gives another
valid bound. Indeed $M_2$ is the exact squared-density residual and
$\mu_r\le\mu_c$ follows from $1/V\le e^{1/V}-1$.
One may take the minimum of any finite list of these valid bounds, then
use the same positive-part formula for its remainder. The higher-moment
version is often considerably sharper than a second-moment estimate alone.

For fixed $T,q,p,\theta$, the allowance $\mathcal B_p(V,\theta)$ is
nonincreasing in $V$: the Gaussian profile, $S_p$, and $E_p$ decrease,
and the reverse-floor transition is a downward jump. The minimum over any
fixed finite comparison family is also nonincreasing. Exact-arithmetic
bisection can therefore bracket that family's passing threshold; this does
not establish optimal actual privacy noise or monotonicity across batch sizes.

\paragraph{Calibration must pay the complete allowance.}
At nominal CLT noise, $g_{\mu_c}(\varepsilon)=\delta$ uses the entire
budget. A positive remainder cannot then be discarded. A corrected
public choice must verify
\begin{equation}\label{eq:finite-corrected-selection}
 V\ge V_{\rm CLT},\qquad \mathcal B_p(V,\theta)\le\delta.
\end{equation}
The condition $\mathcal B_p\le\delta$ alone suffices for finite privacy.
The additional constraint $V\ge V_{\rm CLT}$ defines the upward-correction
policy; it is not a necessary privacy condition. Together these conditions
imply $g_{\mu_c}(\varepsilon)+\mathcal R\le\delta$.
They have an unconditional fallback: take
$V=\max\{V_{\rm CLT},V_\star\}$, $\theta=\varepsilon/2$, and the order
in \eqref{eq:order}. The calibration proof in
Appendix~\ref{sec:privacy} pays the ordinary higher-moment allowance for
every $V\ge V_\star$; the reverse-floor improvement can only reduce it.
At $q=1$, use exact Gaussian composition and $V=T/\mu^2$ instead.
Thus feasibility is supplied by a proved calibration, not assumed of the
nominal CLT variance. Searching public candidates for a smaller passing
variance is legitimate; it is not a proof of minimum feasible privacy noise.

\paragraph{Coverage and utility at the corrected noise.}
On the 28 subsampled grid cells, a finite search with
$p\in\{2,\ldots,32\}$ (also including the fallback order) and
$\theta/\varepsilon\in\{1/20,\ldots,19/20\}$ finds no passing nominal
CLT variance under this sufficient test. Corrected candidates with budget
at most $0.95\delta$ have variance ratios to nominal CLT between about
$1.38$ and $10.40$, and to saved PLD between $1.37$ and $8.80$.
These ranges describe the tested family, not unavoidable privacy penalties.
In particular, failure of this upper-bound screen does not prove failure
of DP; Proposition~\ref{prop:finite-clt-failure} establishes that separately
for one cell. The retained reference budgets were checked independently at 70-digit
precision; the maintained implementation also passed a separate 80-digit
check. These are floating diagnostics, not interval certificates.

For every such public corrected variance, the noise energy remains
$C^2dV/(qN)^2$. Theorem~\ref{thm:main} and Corollary~\ref{cor:rate} therefore
apply to exactly the same compact coefficient and the full cap
$1/[L(1+b)]$. No additional CLT hypothesis or smaller stepsize cap enters
that utility argument. Comparisons at these corrected variances share a
guaranteed target privacy level; their actual privacy losses need not be equal.

\section{A finite failure of the CLT prescription}
\label{sec:finite-clt-failure}
\begin{proposition}\label{prop:finite-clt-failure}
Let $T=1000$, $q=64/50000=4/3125$, and let $\mu>0$ solve
$g_\mu(1)=10^{-5}$. Put
\[
 V=\frac1{\log(1+\mu^2/(Tq^2))},\qquad u=V^{-1/2},\qquad
 Q=\N(0,1),\qquad P=(1-q)Q+q\N(u,1).
\]
For $A=\{\max_{1\le t\le T}Z_t>5\}$,
\[
 P^{\otimes T}(A)-eQ^{\otimes T}(A)>6.9\times10^{-4}>10^{-5}.
\]
Consequently this noise prescription does not give a universal finite
$(1,10^{-5})$ transcript guarantee for the record-query class in
\eqref{eq:algorithm}.
\end{proposition}
\begin{proof}
Write $\overline\Phi(x)=1-\Phi(x)$. For $x>0$, integration by parts gives
the Mills inequalities
\[
 \frac{x\varphi(x)}{1+x^2}<\overline\Phi(x)<\frac{\varphi(x)}x.
\]
Indeed the upper bound follows by replacing $1$ by $z/x$ in the tail
integral; integration by parts gives
$\overline\Phi(x)=\varphi(x)/x-\int_x^\infty\varphi(z)/z^2\,dz$,
which yields the lower bound.
Using $e\varphi(33/8)=\varphi(31/8)$, these imply
\[
 g_{1/4}(1)
 <\varphi(31/8)\left(\frac8{31}-\frac{264}{1153}\right)
 <6.4\times10^{-6}<10^{-5}.
\]
Strict monotonicity of $g_\mu(1)$ gives $\mu>1/4$.
Since $e^{361/100}<1+(1/16)/(Tq^2)$, it follows that $u>1.9$.

Put $\alpha=\overline\Phi(5)$ and
$\beta=(1-q)\alpha+q\overline\Phi(5-u)$.
Independence gives
\[
 Q^{\otimes T}(A)=1-(1-\alpha)^T,\qquad
 P^{\otimes T}(A)=1-(1-\beta)^T.
\]
The following loose constants suffice:
\[
 1.48\times10^{-6}<\varphi(5)<1.49\times10^{-6},\qquad
 \varphi(3.1)>0.00326,\qquad e<2.72.
\]
Thus $\alpha<a_U=1.49\times10^{-6}/5$ and $\beta>b_L$, where
\[
 b_L=(1-q)\frac{5(1.48\times10^{-6})}{26}
       +q\frac{3.1(0.00326)}{1+3.1^2}.
\]
The union bound and $(1-\beta)^T\le e^{-T\beta}$ now give, with $x=Tb_L$,
\[
 P^{\otimes T}(A)-eQ^{\otimes T}(A)
 >1-e^{-x}-2.72Ta_U
 \ge x-\frac{x^2}{2}-2.72Ta_U
 >0.00069175.
\]
All displayed decimal bounds can be checked by rational arithmetic:
use $223/71<\pi<22/7$, hence $2.506<\sqrt{2\pi}<2.508$, and, for
$z>0$ and $n+2>z$,
\[
 \sum_{k=0}^n\frac{z^k}{k!}<e^z<
 \sum_{k=0}^n\frac{z^k}{k!}
 +\frac{z^{n+1}}{(n+1)!}\frac1{1-z/(n+2)}.
\]
Taking $n=80$ verifies the exponential and density bounds above; the final
comparison is rational.

To realize the product experiment, take $N=50000$ slots and a one-dimensional
query equal to $C$ on the sole present record, independent of the state;
all other slots are empty. Its neighbor has every slot empty. The normalized
transcript $Y_t/(C\sqrt V)$ has law $P^{\otimes T}$ or $Q^{\otimes T}$,
respectively, so the displayed event disproves the transcript guarantee.
For larger $d$, the same construction uses the first coordinate.
\end{proof}

This realization concerns the general privacy query class. It does not
satisfy the bounded-below objective assumptions of the utility theorem,
and does not establish failure for every adaptive instance or final model.
Data processing can reduce distinguishability. Direct evaluation at the
inverted value gives $V\approx0.2629208143$, $u\approx1.950237557$ and an
event gap approximately $0.00097140769$; these diagnostics are not needed
for the conservative proof. There is no conflict with the fixed-pass
$G_0$ limit: the finite failure and the limiting excess privacy concern
different points along the parameter family.

\begin{thebibliography}{9}
\bibitem{bu2020deepgdp}
Z. Bu, J. Dong, Q. Long and W. J. Su.
\emph{Deep Learning with Gaussian Differential Privacy}.
Harvard Data Science Review 2(3), 2020.
\href{https://doi.org/10.1162/99608f92.cfc5dd25}{doi:10.1162/99608f92.cfc5dd25}.

\bibitem{clipping}
A. Koloskova, H. Hendrikx and S. U. Stich.
\emph{Revisiting Gradient Clipping: Stochastic Bias and Tight Convergence
Guarantees}. ICML, 2023. \href{https://arxiv.org/abs/2305.01588}{arXiv:2305.01588}.
\bibitem{amap}
E. Monzio Compagnoni, A. Stanghellini, R. Islamov, A. Lucchi and A. Koloskova.
\emph{Adaptive Methods Are Preferable in High Privacy Settings:
An SDE Perspective}. ICLR, 2026.
\href{https://arxiv.org/abs/2603.03226}{arXiv:2603.03226}.
\bibitem{gross}
L. Gross. \emph{Logarithmic Sobolev Inequalities}.
American Journal of Mathematics, 97(4):1061--1083, 1975.
\href{https://doi.org/10.2307/2373688}{doi:10.2307/2373688}.
\bibitem{gdp}
J. Dong, A. Roth and W. J. Su.
\emph{Gaussian Differential Privacy}. Journal of the Royal Statistical
Society: Series B, 84(1):3--37, 2022.
\href{https://arxiv.org/abs/1905.02383}{arXiv:1905.02383}.
\bibitem{abadi}
M. Abadi et al.\ \emph{Deep Learning with Differential Privacy}.
ACM CCS, 2016. \href{https://arxiv.org/abs/1607.00133}{arXiv:1607.00133}.
\bibitem{snm}
O. R\"ais\"a, J. J\"alk\"o and A. Honkela.
\emph{Subsampling is not Magic: Why Large Batch Sizes Work for Differentially
Private Stochastic Optimisation}. ICML, PMLR 235:41959--41981, 2024.
\href{https://proceedings.mlr.press/v235/raisa24a.html}{Published version}.
\end{thebibliography}
\end{document}
